% !TEX TS-program = LuaLaTeX
%%==================================================
%% 这是一个关于bit硕士论文模板的例子
%% 由大眼小蚂蚁根据交大的模板修个得到的
%% Encoding: UTF-8
%% last update: Dec 5th, 2010
%%==================================================

% 字号选项: c5size 五号(默认)
% 字号选项: cs4size 小四
% 双面打印（注意字号设置）
% \documentclass[cs4size, a4paper, cs4size, twoside]{bitmaster-xetex}
% 单面打印（注意字号设置）
\documentclass[cs4size, a4paer, cs4size, twoside, openany]{bitmaster-xetex}


\newboolean{DOIT}
\setboolean{DOIT}{false}
%编译某些只想自己看的内容，编译true，否则false

%% 行距缩放因子(x倍字号)

% 设置图形文件的搜索路径
\graphicspath{{figure/}{figures/}{logo/}{logos/}{graph/}{graphs}}

%%========================================
%% 在bitmaster-xetex.cls中定义的有用命令
%%========================================
% \cndash 中文破折号
% 数学常量
% \me 对数常数e
% \mi 虚数单位i
% \mj 虚数单位j
% \dif 直立的微分算符d为直立体。
% 可伸长的数学箭头、等号
% \myRightarrow{}{}
% \myLeftarrow{}{}
% \myBioarrow{}{}
% \myLongEqual{}{}
% 参考文献
% \upcite{} 上标引用
%%========================================
%调节段落间隔使用
%\inv 压缩间隔
%\smallinv 略微压缩间隔
%\forw 扩大间隔
%\smallforw 略微扩大间隔
\def\inv{\vspace*{-8pt}}
\def\smallinv{\vspace*{-3pt}}
\def\forw{\vspace*{5pt}}
\def\smallforw{\vspace*{3pt}}


\begin{document}

%%%%%%%%%%%%%%%%%%%%%%%%%%%%%%
%% 封面
%%%%%%%%%%%%%%%%%%%%%%%%%%%%%%

% 中文封面内容(关注内容而不是形式)
\classification{TQ028.1}
\UDC{111}

\title{\LARGE{BIT硕士论文\LaTeX 模板}}
\author{大眼小蚂蚁}
\institute{机电学院}
\advisor{教授}
\chairman{离散}
\degree{工学硕士}
\major{兵器科学与技术}
\school{北京理工大学}
\defenddate{2012年12月}
\studentnumber{2120100277}


% 英文封面内容(关注内容而不是表现形式)
\englishtitle{Template for BIT Mater Degree Thesis}
\englishauthor{}
\englishadvisor{Prof.San Zhang}
\englishschool{Beijing Insititute of Technology}
\englishinstitute{Depart of XXX}
\englishdegree{Master}
\englishmajor{Physics}
\englishdate{Dec, 2012}

% 封面,这个命令实现了形式，具体见.cls
\maketitle


% 英文封面
\makeenglishtitle

%打印竖排论文题目
\makeVerticalTitle

% 论文原创性声明和使用授权
\makeDeclareOriginal

%%%%%%%%%%%%%%%%%%%%%%%%%%%%%%
%% 前言
%%%%%%%%%%%%%%%%%%%%%%%%%%%%%%
\frontmatter
\pagenumbering{Roman}

\fancypagestyle{plain}{% 设置开章页页眉页脚风格
  \fancyhf{}%
  \fancyhead[C]{\zihao{5} {\songti 北京理工大学硕士学位论文}}  % 奇数页左页眉
  %\fancyhead[RO]{\small {\it \leftmark}}                   % 奇数页左页眉
  %\fancyhead[LE]{\small {\it \CAST@value@titlemark}} 	 % 偶数页左页眉
  %\fancyhead[RE]{\small {\it 北京理工大学硕士学位论文}}       % 偶数页右页眉
  \fancyfoot[C]{\zihao{5} ~---~{\songti \thepage}~---~} %页脚格式
}

% 摘要
%%==================================================
%% abstract.tex for SJTU Master Thesis
%% based on CASthesis
%% modified by wei.jianwen@gmail.com
%% version: 0.3a
%% Encoding: UTF-8
%% last update: Dec 5th, 2010
%%==================================================

\begin{abstract}

  (摘要是一篇具有独立性和完整性的短文，应概括而扼要地反映出本论文的主要内
  容。包括研究目的、研究方法、研究结果和结论等，特别要突出研究结果和结论。
  中文摘要力求语言精炼准确，硕士学位论文摘要建议500---800字，博士学位论文
  建议1000---1200字。摘要中不可出现参考文献、图、表、化学结构式、非公知公
  用的符号和术语。英文摘要与中文摘要的内容应一致。）

  \keywords{北京理工 \quad 联盟论坛 \quad 爱国荣校}
\end{abstract}

\begin{englishabstract}

  this is english abstract

  \englishkeywords{BIT, master thesis, XeTeX/LaTeX template}
\end{englishabstract}


% 目录
\tableofcontents
% 表格索引
%\listoftables
% 插图索引
%\listoffigures

%\addcontentsline{toc}{chapter}{\listfigurename} %将表格索引加入全文目录
%\addcontentsline{toc}{chapter}{\listtablename}  %将图索引加入全文目录

% 主要符号、缩略词对照表
% %%==================================================
%% symbol.tex for SJTU Master Thesis
%% based on CASthesis
%% modified by wei.jianwen@gmail.com
%% version: 0.3a
%% Encoding: UTF-8
%% last update: Dec 5th, 2010
%%==================================================

\chapter{主要符号对照表}
\label{chap:symb}
\begin{tabular}{ll}

 \hspace{2em}$\epsilon$       & \hspace{5em}介电常数 \\
 \hspace{2em}$\mu$ \qquad     & \hspace{5em}磁导率 \\
  \hspace{2em}$\epsilon$       & \hspace{5em}介电常数 \\
 \hspace{2em}$\mu$ \qquad     & \hspace{5em}磁导率 \\
 \hspace{2em}$\epsilon$       & \hspace{5em}介电常数 \\
 \hspace{2em}$\mu$ \qquad     & \hspace{5em}磁导率 \\
 \hspace{2em}$\epsilon$       & \hspace{5em}介电常数 \\
 \hspace{2em}$\mu$ \qquad     & \hspace{5em}磁导率 \\


\end{tabular}


%%%%%%%%%%%%%%%%%%%%%%%%%%%%%%
%% 正文
%%%%%%%%%%%%%%%%%%%%%%%%%%%%%%
\mainmatter

\fancypagestyle{plain}{% 设置开章页页眉页脚风格
  \fancyhf{}%
  \fancyhead[C]{\zihao{5} {\songti 北京理工大学硕士学位论文}}  % 奇数页左页眉
  %\fancyhead[RO]{\small {\it \leftmark}}                   % 奇数页左页眉
  %\fancyhead[LE]{\small {\it \CAST@value@titlemark}} 	 % 偶数页左页眉
  %\fancyhead[RE]{\small {\it 北京理工大学硕士学位论文}}       % 偶数页右页眉
  \fancyfoot[C]{\zihao{5} ~---~{\songti \thepage}~---~} %页脚格式
}


%% 各章正文内容
%%==================================================
%% chapter01.tex for BIT Master Thesis
%% based on CASthesis
%% modified by tingzhen.li@gmail.com
%% version: 0.1
%% Encoding: UTF-8
%% last update: Dec, 2012
%%==================================================

\chapter{这是什么}
\label{chap:what}

这是北京理工大学硕士学位学位论文~\LaTeX~模板，其由上海交大版本
\footnote{网页见\url{http://blog.sina.com.cn/s/blog_5e16f1770100nsqy.html}}修改得到，
修改依据《北京理工大学博士、硕士学位论文撰写规范》。

\section{模板说明}
\label{sec:fastguide}

\subsection{模板特性}
\label{sec:features}

这个模板基于~CASthesis--0.1j~文档类，中文解决方案是~\XeTeX/\LaTeX~。参考
文献建议使用~BibTeX~管理，可以生成符合国标~GBT7714~风格的参考文献列表。模
板在~Windows~和~Linux~下测试通过，更详细的系统要求请参考
\ref{sec:requirements}。

模板的外观表现和功能都放
在~bitmaster-xetex.cls~和~bitmaster-xetex.cfg~中，在对外观进行细微调整
时，只需要更新这两个文件，不需要对.tex源文件做修改。这也给模板更新带来了
极大方便。

最后，给出一个列表，罗列一下这个模板的功能要点：

\begin{itemize}
\item \inv  使用~\XeTeX~引擎处理中文，{\color{red} 我在ubuntu上安装了
  texlive2011和windows中文字体，具体方法可参见\url{http://blog.163.com/warrior511@126/blog/static/16798651220114984811658/}}；
\item \inv 包含中文字符的源文件(.tex, .bib, .cfg)，编码都使用UTF-8；
\item \inv 使用~BibTeX~管理参考文献。参考文献表现形式(格式)受~.bst~控制，方便
  在不同风格间切换，目前生成的列表符合国标GBT7714要求；
\item \inv 可以直接插入EPS/PDF/JPG/PNG格式的图像，并且\emph{不需要}~bounding box~文件(.bb)。
\item \inv 模板的格式受~bitmater-xetex.cls~和~bitmaster-xetex.cfg~控制，方便模
  板更新和模板修改。
\end{itemize}

\subsection{系统要求}
\label{sec:requirements}

要使用这个模板协助你完成研究生学位论文的创作，下面的条件必须满足：

\begin{itemize}
\item \inv 操作系统字体目录中有中文字体（adobe或windows字体均可）；
\item \inv \TeX~系统有~\XeTeX~引擎；
\item \inv \TeX~系统有~ctex~宏包；
\item \inv \TeX~系统的fontspec宏包(fontspec.sty)足够新；
\item \inv 你有使用~\LaTeX~的经验。
\end{itemize}

 
\subsection{模板文件布局}
\label{sec:layout}

\begin{lstlisting}[basicstyle=\small\ttfamily,caption={模板文件布局},label=layout,float,numbers=none]
  ├─ diss.tex
  ├── README.pdf
  ├── bitmaster-xetex.cfg
  ├── bitmaster-xetex.cls
  ├── body
  │   ├── abstract.tex
  │   ├── app1.tex
  │   ├── app2.tex
  │   ├── chapter01.tex
  │   ├── chapter02.tex
  │   ├── conclusion.tex
  │   ├── projects.tex
  │   ├── pub.tex
  │   ├── resume.tex
  │   ├── symbol.tex
  │   └── thanks.tex
  ├── figures
  │   └── chap2
  ├── GBT7714-2005NLang.bst
  ├── reference
  │   ├── chap1.bib
  │   └── chap2.bib
  ├── run.sh
  
      
\end{lstlisting}

你拿到手的模板文件大致会包含代码\ref{layout}所列的文件，乍看起来还是挺令
人头大的。并且，这还是“干净”的时候，等到真正开始处理的时候，会冒出相当多
的“中间文件”，这又会使情况变得更糟糕。所以，有必要对这些文件做一些简要说
明。看完这部分以后，你应该发现，其实你要关心的文件类型并没有那么多。

\subsubsection{格式控制文件}
\label{sec:format}

格式控制文件控制着论文的表现形式，包括以下几个文件：
bitmaster-xetex.cfg, bitmaster-xetex.cls~和~GBT7714-2005NLang.bst。
其中，``.cfg''和``.cls''控制论文主体格式，``.bst''控制参考文献条目的格式，

一般用户最好``忽略''格式控制文件的存在，不要去碰它们。
有其他格式需要，欢迎到板上发贴。
对于因为擅自更改格式控制文件出现的问题，我一概不负责。{\large\smiley}

\subsubsection{主控文件~demo.tex}
\label{sec:demotex}

主控文件~demo.tex~的作用就是将你分散在多个文件中的内容``整合''成一篇完整的论文。
使用这个模板撰写学位论文时，你的学位论文内容和素材会被``拆散''到各个文件中：
譬如各章正文、各个附录、各章参考文献等等。
在~demo.tex~中通过``include''命令将论文的各个部分包含进来，从而形成一篇结构完成的论文。
封面页中的论文标题、作者等中英文信息，也是在~demo.tex~中填写。
部分可能会频繁修改的设置，譬如行间距、图片文件目录等，我也放在了demo.tex中。
你也可以在demo.tex中按照自己的需要引入一些的宏包
\footnote{我对宏包的态度是：只有当你需要在文档中使用那个宏包时，才需要在导言区中用~usepackage~引入该宏包。如若不然，通过usepackage引入一大堆不被用到的宏包，必然是一场灾难。由于一开始没有一致的设计目标，\LaTeX~的各宏包几乎都是独立发展起来的，因重定义命令导致的宏包冲突屡见不鲜。}。

大致而言，在~demo.tex~中，大家只要留意把``章''一级的内容，以及各章参考文
献内容包含进来就可以了。需要注意，处理文档时所有的操作命令
{}\cndash{}xelatex, bibtex等，都是作用在~demo.tex~上，而\emph{不是}后面这
些``分散''的文件，请参考\ref{sec:process}小节。

\subsubsection{论文主体文件夹body}
\label{sec:thesisbody}

这一部分是论文的主体，是以``章''为单位划分的。

正文前部分(frontmatter)：中英文摘要(abstract.tex)。其他部分，诸如中英文封
面、授权信息等，都是根据~demo.tex~所填的信息``画''好了，不单独弄成文件。

正文部分(mainmatter)：自然就是各章内容~chapter\emph{xxx}.tex~了，这部分无法自动生成{\LARGE\Smiley}

正文后的部分(backmatter)：附录(app\emph{xx}.tex)；致谢(thuanks.tex)；攻读
学位论文期间发表的学术论文目录(pub.tex)；个人简历(resume.tex)。参考文献列
表是``生成''的，也不作为一个单独的文件。另外，学校的硕士研究生学位论文模
板中，也没有要求加入个人建立，所以我没有在~demo.tex~中引入resume.tex。

\subsubsection{图片文件夹~figures}
\label{sec:figuresdir}

figures~文件夹放置了需要插入文档中的图片文件(PNG/JPG/PDF/EPS)，建议按章再
划分子目录。

\subsubsection{参考文献数据库文件夹~reference}
\label{sec:bibdir}

reference~文件夹放置的是各章``可能''会被引用的参考文献文件。参考文献的元
数据，例如作者、文献名称、年限、出版地等，会以一定的格式记录在纯文本文
件.bib中。最终的参考文献列表是BibTeX处理.bib后得到的，名为~demo.bbl。将参
考文献按章划分的一个好处是，可以在各章后生成独立的参考文献，不过，现在看
来没有这个必要。关于参考文献的管理，可以进一步参考第\ref{chap:example}章
中的例子。

\subsection{如何使用}
\label{sec:process}

模板使用~\XeTeX~引擎提供的~xelatex~的命令处理，作用于“主控文档”demo.tex。
并且，可以省略扩展名。完整的处理流程是：
\begin{enumerate}
\item[] \inv xelatex -no-pdf -{}-interaction=nonstopmode demo
\item[] \inv bibtex demo 
\item[] \inv xelatex -no-pdf -{}-interaction=nonstopmode demo 
\item[] \inv xelatex -{}-interaction=nonstopmode demo 
\end{enumerate}

运行bibtex的时候会提示一些错误，猜测是~{{\sc Bib}\TeX}~对UTF-8支持不充
分，一般不影响最终结果。留意因为拼写错误导致的``找不到文献错误''即可。

加入~\verb|--interaction=nonstopmode|~参数是不让错误打断编译过程。
\XeTeX~仍存在一些宏包兼容性问题，所以会产生一些莫名其妙的错误(通常是重定
义错误)，而这些错误通常不会影响最终的编译结果\footnote{xunicode宏包就很蛋
  痛地重定义了几个数学符号，还有诸如$\backslash$nobreakspace命令}。为方便
使用，我把处理过程写到了run.sh(for Linux)和run.bat(for Windows)批处理文件
中。更规范的应该是使用Makefile，可惜笔者功力不够，希望哪位高人把好用的
Makefile补上。

基本处理流程就是这样，一些~\LaTeX~排版的小例子可以参考第二章。

\section{硕士学位论文格式的一些说明}
\label{sec:thesisformat}


所有关于研究生学位论文模板的要求，我参考了《北京理工大学博士、硕士学位论
  文撰写规范》和研究生院提供的相关参考模板。可惜的是，很多地方参考模板同
规范相违背。例如，规范中指出，目录一级节标题采用小四、加粗，而在模板中并未
加粗。对于这种问题，基本上以大多数同学采用的方式保持一致，没有严格按照规
范进行。

这个模板是为``单面打印''准备的，但也支持``双面打印''。你可以在~demo.tex~
中设定文档类的语句中进行相应修改：

\begin{quote}
  {\small\verb+\documentclass[cs4size, a4paer, cs4size, oneside, openany]{bitmaster-xetex}+}
\end{quote}


关于页眉页脚，按照BIT要求：页眉为``北京理工大学XX学位论文''，XX表示博士或
硕士，宋体、5号，居中排列；页眉从中文摘要开始标注，论文页眉奇偶页相同。页
码从第1章（绪论）开始按阿拉伯数字（1，2，3……）连续编排，之前的部分(中文摘
要、Abstract、目录等)用大写罗马数字（Ⅰ，Ⅱ，Ⅲ……）单独编排。

研究生院要求参考文献必须符合~GBT7714~风格，学校明确提出使用这个标准而不是
自己拍脑袋想出别的做法，应该算是谢天谢地了。使用这个模板，结合BibTeX，可
以很方便地生成符合GB标准的参考文献列表。

\section{bit-master-thesis模板类简介}
\label{sec:bit.cls}
论文模板主要在bit-master-thesis.cls文件中进行定义，现对其进行简单介绍。
\subsection{页面设置}
页边距设置如下，{\color{red}好像存在问题，需进一步解决}：
\begin{lstlisting}
\usepackage[top=3.5cm,headheight=25mm,headsep=3mm,footskip=8mm,bottom=2.5cm,left=2.7cm,right=2.7cm]{geometry}
\end{lstlisting}

行距离设置，按照要求，应该为22榜，如下设置，效果基本相同：
\begin{lstlisting}
\RequirePackage{setspace}
\setstretch{1.4}
\end{lstlisting}

\subsection{章节格式与目录}
严格按照规范，采用如下代码实现：
\begin{lstlisting}
%% 设置章节格式
\CTEXsetup[number={\arabic{chapter}},name={第,章},
            nameformat={\bfseries\heiti\centering\zihao{3}},
            titleformat={\bfseries\heiti\zihao{3}},
            afterskip={30pt}]{chapter}
\CTEXsetup[nameformat={\bfseries\heiti\zihao{4}},
            titleformat={\bfseries\heiti\zihao{4}}]{section}
\CTEXsetup[nameformat={\bfseries\heiti\zihao{-4}},
            titleformat={\bfseries\heiti\zihao{-4}}]{subsection}
\CTEXsetup[nameformat={\bfseries\zihao{-4}},
            titleformat={\bfseries\zihao{-4}}]{subsubsection}
\CTEXsetup[format={\Large\bfseries}]{section}
\CTEXsetup[beforeskip={10pt}]{chapter}

%% 用\textsf{titletoc}设定目录格式。
\RequirePackage{titletoc}
\titlecontents{chapter}[0pt]{\vspace{0.25\baselineskip} \songti \zihao{4}}
    {\thecontentslabel\quad}{}
    {\hspace{.5em}\titlerule*{.}\contentspage}
\titlecontents{section}[2em]{\songti \zihao{-4}}
    {\thecontentslabel\quad}{}
    {\hspace{.5em}\titlerule*{.}\contentspage}
\titlecontents{subsection}[4em]{\songti \zihao{-4}}
    {\thecontentslabel\quad}{}
    {\hspace{.5em}\titlerule*{.}\contentspage}
\end{lstlisting}

\subsection{封面设计}
这里我们为封面设计提供了众多命令，以中文封面为例：
\begin{lstlisting}
%%%%中文标题页的可用命令
\newcommand\classification[1]{\def\CAST@value@classification{#1}}
\newcommand\studentnumber[1]{\def\CAST@value@studentnumber{#1}}
\newcommand\confidential[1]{\def\CAST@value@confidential{#1}}
\newcommand\UDC[1]{\def\CAST@value@UDC{#1}}
\newcommand\serialnumber[1]{\def\CAST@value@serialnumber{#1}}
\newcommand\school[1]{\def\CAST@value@school{#1}}
\newcommand\degree[1]{\def\CAST@value@degree{#1}}
\renewcommand\title[2][\CAST@value@title]{%
  \def\CAST@value@title{#2}
  \def\CAST@value@titlemark{\MakeUppercase{#1}}}
\renewcommand\author[1]{\def\CAST@value@author{#1}}
\newcommand\advisor[1]{\def\CAST@value@advisor{#1}}
\newcommand\advisorinstitute[1]{\def\CAST@value@advisorinstitute{#1}}
\newcommand\major[1]{\def\CAST@value@major{#1}}
\newcommand\submitdate[1]{\def\CAST@value@submitdate{#1}}
\newcommand\defenddate[1]{\def\CAST@value@defenddate{#1}}
\newcommand\institute[1]{\def\CAST@value@institute{#1}}
\newcommand\chairman[1]{\def\CAST@value@chairman{#1}}
\end{lstlisting}

使用这些命令，即可在主控文件中设置自己的封面，例如本文档在demo.tex中如下
设置：
\begin{lstlisting}
\classification{TQ028.1}
\UDC{111}
\title{\LARGE{BIT硕士论文\LaTeX 模板}}
\author{大眼小蚂蚁}
\institute{机电学院}
\advisor{教授}
\chairman{离散}
\degree{工学硕士}
\major{兵器科学与技术}
\school{北京理工大学}
\defenddate{2012年12月}
\studentnumber{2120100277}
\end{lstlisting}

这些变量设置好之后，\\使用\verb+\maketitle+产生封面的第一二页；\\使用
\verb+\makeenglishtitle+产生英文标题页；\\使用\verb+\makeVerticalTitle+产
生竖着排放的标题页;\\使用\verb+\makeDeclareOriginal+产生声明页。


%%==================================================
%% chapter02.tex for BIT Master Thesis
%% based on CASthesis
%% modified by wei.jianwen@gmail.com
%% modified by 大眼小蚂蚁
%% version: 0.3a
%% Encoding: UTF-8
%% last update: Dec 5th, 2010
%%==================================================

% \bibliographystyle{bit2} %[此处用于每章都生产参考文献]

\chapter{一些~\LaTeX~排版的例子}
\label{chap:example}

\section{数学排版的例子}
\label{sec:matheq}

\subsection{公式排版}
\label{sec:eqformat}

这里有举一个长公式排版的例子，来自\href{http://www.tex.ac.uk/tex-archive/info/math/voss/mathmode/Mathmode.pdf}{《Math mode》}
\footnote{《Math mode》中的例子实在太丰富了，而且每一个都很精彩实用。我本想抄几个例子上来做个``山寨版''，实在没有必要，大家还是去看``原著''吧。}：

\begin {multline}
  \frac {1}{2}\Delta (f_{ij}f^{ij})=
  2\left (\sum _{i<j}\chi _{ij}(\sigma _{i}-
    \sigma _{j}) ^{2}+ f^{ij}\nabla _{j}\nabla _{i}(\Delta f)+\right .\\
  \left .+\nabla _{k}f_{ij}\nabla ^{k}f^{ij}+
    f^{ij}f^{k}\left [2\nabla _{i}R_{jk}-
      \nabla _{k}R_{ij}\right ]\vphantom {\sum _{i<j}}\right )
\end{multline}

\subsubsection{一个四级标题}
\label{sec:depth4}

这是全文唯一的一个四级标题。在这部分中将演示可伸长符号（箭头、等号的例子）的例子，以及如何在可伸长的符号上标注。在\href{http://zhou63.ahut.edu.cn/latex/ctexfaq.pdf}{《CTeX常见问题集》}中也由类似的介绍。
首先需要在~diss.tex~导言区引入如下的内容：

\begin{lstlisting}[language={TeX}, caption={插入导言区的内容}]
  \makeatletter
  \def\ExtendSymbol#1#2#3#4#5{\ext@arrow 0099{\arrowfill@#1#2#3}{#4}{#5}}
  \def\RightExtendSymbol#1#2#3#4#5{\ext@arrow 0359{\arrowfill@#1#2#3}{#4}{#5}}
  \def\LeftExtendSymbol#1#2#3#4#5{\ext@arrow 6095{\arrowfill@#1#2#3}{#4}{#5}}
  \makeatother
  
  \newcommand\myRightarrow[2][]{\RightExtendSymbol{=}{=}{\Rightarrow}{#1}{#2}}
  \newcommand\myLeftarrow[2][]{\LeftExtendSymbol{\Leftarrow}{=}{=}{#1}{#2}}
  \newcommand\myBioarrow[2][]{\ExtendSymbol{\Leftarrow}{=}{\Rightarrow}{#1}{#2}}
  \newcommand\myLongEqual[2][]{\ExtendSymbol{=}{=}{=}{#1}{#2}}
\end{lstlisting}

然后，在正文插入如代码\ref{extend}所示的内容。效果如下：

\begin{lstlisting}[language={TeX}, caption={可伸长的符号}, label=extend]
  \begin{eqnarray}
    f(x) & \myBioarrow{A=B}  & B \\
    & \myLongEqual{A=B} & B \\
    & \myLeftarrow[A=B^2]{B=A^2} & B \nonumber \\
    & \myRightarrow{B^2=A^2} & B
  \end{eqnarray}
\end{lstlisting}

\begin{displaymath}
    A \xleftarrow{n=0} B \xrightarrow[LongLongLongLong]{n>0} C 
\end{displaymath}

\begin{eqnarray}
  f(x) & \myBioarrow{A=B}  & B \\
  & \myLongEqual{A=B} & B \\
  & \myLeftarrow[A=B^2]{B=A^2} & B \nonumber \\
  & \myRightarrow{B^2=A^2} & B
\end{eqnarray}

又如：

\begin{align}
  \label{eq:none}
  & I(X_3;X_4)-I(X_3;X_4|X_1)-I(X_3;X_4|X_2) \nonumber \\
  \myLongEqual{a)}\, & [I(X_3;X_4)-I(X_3;X_4|X_1)]-I(X_3;X_4|\tilde{X}_2) \\
  \myLongEqual[\rule{0.28cm}{0cm}]{}\, & I(X_1;X_3;X_4)-I(X_3;X_4|\tilde{X}_2)
\end{align}


\subsection{定理环境}

在~CASthesis.cfg~中定义了丰富的定理环境
algo(算法)，thm(定理)，lem(引理)，prop(命题)，cor(推论)，defn(定义)，conj(猜想)，exmp(例)，rem(注)，case(情形)，
bthm(断言定理)，blem(断言引理)，bprop(断言命题)，bcor(断言推论)。
amsmath还提供了一个proof(证明)的环境。
这里举一个``定理''和``证明''的例子。
\begin{thm}[留数定理]
\label{thm:res}
  假设$U$是复平面上的一个单连通开子集，$a_1,\ldots,a_n$是复平面上有限个点，$f$是定义在$U\backslash \{a_1,\ldots,a_n\}$上的全纯函数，
  如果$\gamma$是一条把$a_1,\ldots,a_n$包围起来的可求长曲线，但不经过任何一个$a_k$，并且其起点与终点重合，那么：

  \begin{equation}
    \label{eq:res}
    \ointop_{\gamma}f(z)\,\mathrm{d}z = 2\uppi\mathbf{i}\sum^n_{k=1}\mathrm{I}(\gamma,a_k)\mathrm{Res}(f,a_k)
  \end{equation}

  如果$\gamma$是若尔当曲线，那么$\mathrm{I}(\gamma, a_k)=1$，因此：

  \begin{equation}
    \label{eq:resthm}
    \ointop_{\gamma}f(z)\,\mathrm{d}z = 2\uppi\mathbf{i}\sum^n_{k=1}\mathrm{Res}(f,a_k)
  \end{equation}

      % \oint_\gamma f(z)\, dz = 2\pi i \sum_{k=1}^n \mathrm{Res}(f, a_k ). 

  在这里，$\mathrm{Res}(f, a_k)$表示$f$在点$a_k$的留数，$\mathrm{I}(\gamma,a_k)$表示$\gamma$关于点$a_k$的卷绕数。
  卷绕数是一个整数，它描述了曲线$\gamma$绕过点$a_k$的次数。如果$\gamma$依逆时针方向绕着$a_k$移动，卷绕数就是一个正数，
  如果$\gamma$根本不绕过$a_k$，卷绕数就是零。

  定理\ref{thm:res}的证明。
  
  \begin{proof}
    首先，由……

    其次，……

    所以……
  \end{proof}
\end{thm}

上面的公式例子中，有一些细节希望大家注意。微分号d应该使用``直立体''，也就是用mathrm包围起来。
并且，微分号和被积函数之间应该有一段小间隔，可以插入\verb+\,+得到。
斜体的$d$通常只作为一般变量。
i,j作为虚数单位时，也应该使用``直立体''，为了明显，还加上了粗体，例如\verb+\mathbf{i}+。斜体$i,j$通常用作表示``序号''。
其他字母在表示常量时，也推荐使用``直立体''，譬如，圆周率$\uppi$（需要upgreek宏包），自然对数的底$\mathrm{e}$。
不过，我个人觉得斜体的$e$和$\pi$很潇洒，在不至于引起混淆的情况下，我也用这两个字母的斜体表示对应的常量。


\section{向文档中插入图像}
\label{sec:insertimage}

\subsection{支持的图片格式}
\label{sec:imageformat}

\XeTeX~可以很方便地插入~PDF、EPS、PNG、JPG~格式的图片。

插入PNG/JPG的例子如\ref{fig:SRR}所示。
这两个水平并列放置的图共享一个``图标题''(table caption)，没有各自的小标题。

\begin{figure}[!htp]
  \centering
  \includegraphics[width=0.3\textwidth]{chap2/testpng}
  \hspace{1cm}
  \includegraphics[width=0.3\textwidth]{chap2/testjpg}
  \bicaption[fig:SRR]{这里将出现在插图索引中}{中文题图}{Fig}{English caption}
\end{figure}

这里还有插入eps图像和pdf图像的例子，如图\ref{fig:pdfeps}。这里将EPS和PDF图片作为子图插入，每个子图有自己的小标题。并列子图的功能是使用~subfigure~宏包提供的。

\begin{figure}
  \centering
  \subfigure[EPS Figure]{
    \label{fig:epspdf:a} %% label for first subfigure
    \includegraphics[width=0.3\textwidth]{chap2/testeps}}
  \hspace{1in}
  \subfigure[PDF Figure]{
    \label{fig:epspdf:b} %% label for second subfigure
    \includegraphics[angle=-90,origin=br,width=0.3\textwidth]{chap2/testpdf.pdf}}
  \bicaption[fig:pdfeps]{插入eps图像和pdf图像}{插入eps和pdf的例子}{Fig}{An EPS and PDF demo}
\end{figure}

更多关于~\LaTeX~插图的例子可以参考\href{http://www.cs.duke.edu/~junhu/Graphics3.pdf}{《~\LaTeX~插图指南》}。

\subsection{长标题的换行}
\label{sec:longcaption}

图\ref{fig:longcaptionbad}和图\ref{fig:longcaptiongood}都有比较长图标题，通过对比发现，图\ref{fig:longcaptiongood}的换行效果更好一些。
其中使用了minipage环境来限制整个浮动题的宽度。

\begin{figure}[!htp]
 \centering
 \includegraphics[angle=-90,origin=br,width=4cm]{chap2/testpdf.pdf}
 \bicaption[fig:longcaptionbad]{这里将出现在插图索引}{BIT是我国历史最悠久的高等学府之一，是教育部直属、工信部共建的全国重点大学，985，211.}{Fig}{Joomla! is one of the most powerful Open Source Content Management Systems on the planet.}
\end{figure}


  \begin{figure}[!hbp]
    \centering
    \begin{minipage}[b]{0.6\textwidth}
      \captionstyle{\centering}
      \centering
      \includegraphics[angle=-90,origin=br,width=4cm]{chap2/testpdf.pdf}
      \bicaption[fig:longcaptiongood]{这里将出现在插图索引}{BIT是我国历史最悠久的高等学府之一，是教育部直属、工信部共建的全国重点大学，985，211.}{Fig}{Joomla! is one of the most powerful Open Source Content Management Systems on the planet.}
    \end{minipage}     
  \end{figure}

  
\section{表格的例子}
\label{sec:tab}

这一节给出的是一些表格的例子，如表\ref{tab:firstone}所示。

\begin{table}[!hpb]
  \centering
  \bicaption[tab:firstone]{指向一个表格的表目录索引}{一个颇为标准的三线表格\footnotemark[1]}{Table}{A Table}
  \begin{tabular}{@{}llr@{}} \toprule
    \multicolumn{2}{c}{Item} \\ \cmidrule(r){1-2}
    Animal & Description & Price (\$)\\ \midrule
    Gnat & per gram & 13.65 \\
    & each & 0.01 \\
    Gnu & stuffed & 92.50 \\
    Emu & stuffed & 33.33 \\
    Armadillo & frozen & 8.99 \\ \bottomrule
  \end{tabular}
\end{table}
\footnotetext[1]{这个例子来自\href{http://www.ctan.org/tex-archive/macros/latex/contrib/booktabs/booktabs.pdf}{《Publication quality tables in LATEX》}(booktabs宏包的文档)。这也是一个在表格中使用脚注的例子，请留意与threeparttable实现的效果有何不同。}

下面一个是一个更复杂的表格，用threeparttable实现带有脚注的表格，如表\ref{tab:footnote}。

\begin{table}[!htpb]
  \bicaption[tab:footnote]{出现在表目录的标题}{一个带有脚注的表格的例子}{Table}{A Table with footnotes}
  \centering
  \begin{threeparttable}[b]
     \begin{tabular}{ccd{4}cccc}
      \toprule
      \multirow{2}{6mm}{total}&\multicolumn{2}{c}{20\tnote{1}} & \multicolumn{2}{c}{40} &  \multicolumn{2}{c}{60}\\
      \cmidrule(lr){2-3}\cmidrule(lr){4-5}\cmidrule(lr){6-7}
      &www & k & www & k & www & k \\
      \midrule
      &$\underset{(2.12)}{4.22}$ & 120.0140\tnote{2} & 333.15 & 0.0411 & 444.99 & 0.1387 \\
      &168.6123 & 10.86 & 255.37 & 0.0353 & 376.14 & 0.1058 \\
      &6.761    & 0.007 & 235.37 & 0.0267 & 348.66 & 0.1010 \\
      \bottomrule
    \end{tabular}
    \begin{tablenotes}
    \item [1] the first note.% or \item [a]
    \item [2] the second note.% or \item [b]
    \end{tablenotes}
  \end{threeparttable}
\end{table}

\section{参考文献管理}

参考文献的管理是这个学位论文模板又一个好玩的地方。

\subsection{将参考文献的内容与表现分离}

这个论文模板使用BibTeX处理参考文献，这又是一个``内容''与``表现形式''分离的极好例子
\footnote{当然，你也可以手动编参考文献item，直接插入文档中。但是，有BibTeX帮助，我觉得没有人想用这种麻烦的方法，所以就在脚注中说明了。}。
参考文献的``内容''就是reference文件夹下的chap\textit{xx}.bib，参考文献的元数据(名称、作者、出处等)以一定的格式保存在这些纯文本文件中。
.bib文件也可以理解为参考文献的``数据库''，正文中所有引用的参考文件条目都会从这些文件中``析出''。
控制参考文献条目``表现形式''(格式)的是.bst文件。.bst文件定义了参考文献风格，使用不同的参考文献风格能将同一个参考文献条目输出成不同的格式。
当然，一个文档只能使用一个参考文献风格。
按照教务处的要求，本模板使用的是国标GBT7714风格的参考文献。

BibTeX的工作过程是这样的：
BibTeX读取.aux(第一次运行latex得到的)看看你引用了什么参考文献条目，
然后到.bib中找相关条目的信息，
最后根据.bst的格式要求将参考文献条目格式化输出，写到.bbl文件中。
在运行latex将.bbl插入文档之前，你可以用文本编辑器打开它，做一些小的修改。
你会发现，.bbl的格式和你自己手动写item很相似，它已经被赋予了一定的``表现形式''。

.bib数据库中的参考文献条目可以手动编写，也可以在google的学术搜索中找到。
各大数据库\footnote{应该说是国际知名数据库，譬如~SCOPUS, IEEE, OSA等，国内数据库在搜索、导出方面一直是差得一塌糊涂。}也支持将参考文献信息导出为.bib，
省时省力。
以Google学术搜索为例：进入\url{http://scholar.google.cn}，在``学术搜索设置''中，将``文献管理软件''设为``显示导入BibTeX''的连接，保存退出。
然后学术搜索找到文献下会有``导出到BibTeX''连接，点击后Firefox会打开新的标签页，出现类似代码\ref{googlescholar}所示的内容
\footnote{展示这些.bib条目使用了listings宏包，因为listings宏包协调中文的能力很糟糕，所以读者在查看模板的这部分源代码时会看到一些非常麻烦的东西。并且，直接将源代码的这部分内容复制到.bib中可能还会出错。我的建议是：这部分内容留意PDF就足够了。}。
请注意，这个条目离``规范''还有一些距离。

  \begin{lstlisting}[caption={从Google Scholar找到的，但并不规范的.bib条目}, label=googlescholar, float, escapeinside="", numbers=none]
    @phdthesis{"白2008信用风险传染模型和信用衍生品的定价",
      title={{"信用风险传染模型和信用衍生品的定价"}},
      author={"白云芬"},
      year={2008},
      school={"上海交通大学"}
    } 
  \end{lstlisting}

  上面的.bib条目的``名字''\cndash{}``白2008信用风险传染模型和信用衍生品的定价''，包含~ASCII~以外的字符，BibTeX无法处理；
  条目还缺少了address域，这样编译出来的结果会出现``地址不详''；
  并且，条目还缺少language域，BibTeX需要language域来判断是否是中文参考文献。
  将上面的条目修正(改英文名、增加address和language域)，复制到本地的.bib文件中就可以了。
  显然，这里描述的是参考文献的内容，而不是表现形式。

  \begin{lstlisting}[caption={一个符合规范的.bib条目}, label=itemok, float, escapeinside="", numbers=none]
    @phdthesis{bai2008,
      title={{"信用风险传染模型和信用衍生品的定价"}},
      author={"白云芬"},
      year={2008},
      language={zh},
      address={"上海"},
      school={"上海交通大学"}
    } 
  \end{lstlisting}

由于中英文参考文献处理起来有差异，所以需要在参考文献中标注是否是中文文献。
确切地说，BibTeX并不具有区分中英文参考文献的``智能''，这种智慧的来源是.bst文——它定义了处理参考文献的规则。
GBT7714-2005NLang.bst中规定：.bib中的条目，如果条目的``language''域非空，就被认为是中文文献，否则被认为是英文文献。
例如，刚才的文献，就会被认为是中文参考文献，采取一些针对中文的处理方式。

最后，这个条目被bibtex处理后，赋予了一定的``表现形式''，在.bbl文件中以下面的样子出现。
你还可以对它进行小的修改，这是一种很折磨人的终极修改方法。
再次运行latex之后，它将被插入到文档中。

\begin{lstlisting}[caption={.bbl中被格式化之后的条目}, escapeinside="", numbers=none]
\bibitem["白云芬(2008)"]{bai2008}
  \textsc{"白云芬"}.
  \newblock {"信用风险传染模型和信用衍生品的定价"}[D].
  \newblock "上海: 上海交通大学, 2008."
\end{lstlisting}

再罗嗦两句，
.bst文件书写起来非常繁杂\footnote{可以参考《Tame The BeaST》。}，书写符合GBT7714标准的.bst文件更是一项浩大的工程。
因此，当大家为漂亮、标准的参考文献列表感到满意时，应该对
GBT7714-2005NLang.bst的作者充满谢意。{\color{red} 同时强调一下，由于学校
  规定参考文献采用5号字体，所以在使用此规范时，需要在前面加上字号限制：
\begin{lstlisting}
{
\zihao{5}
\bibliography{reference/chap1,reference/chap2}
}
\end{lstlisting}}

再多说两句，.bib是“参考文献的内容”，而控制参考文献表现(格式)的是.bst文件，本模板附带的是GBT7714-2005NLang.bst。

\subsection{在正文中引用参考文献}

参考文献可以分章节管理，只需要在主文件中的参考文献中都包含进去就可以，如\verb+\bibliography{chap1,chap2,...}+。

正文中引用参考文献时，用\verb+\upcite{key1,key2,key3...}+可以产生“上标引用的参考文献”，
如\upcite{Meta_CN,chen2007act,DPMG}。
使用\verb+\cite{key1,key2,key3...}+则可以产生水平引用的参考文献，例如\cite{JohnD,zhubajie,IEEE-1363}。
请看下面的例子，将会穿插使用水平的和上标的参考文献：关于书的\cite{Meta_CN,JohnD,IEEE-1363}，关于期刊的\upcite{chen2007act,chen2007ewi}，
会议论文\cite{DPMG,kocher99,cnproceed}，
硕士学位论文\cite{zhubajie,metamori2004}，博士学位论文\upcite{shaheshang,FistSystem01,bai2008}，标准文件\cite{IEEE-1363}，技术报告\upcite{NPB2}，电子文献\cite{xiaoyu2001, CHRISTINE1998}。

最后总结一些注意事项：
\begin{itemize}
\item  \inv 参考文献只有在正文中被引用了，才会在最后的参考文献列表中出现；
\item \inv 参考文献``数据库文件''.bib是纯文本文件，请使用~UTF-8~编码，不要使用~GBK~编码；
\item \inv 参考文献条目中通过~language~域是否为空判断是否是中文文献；
\item \inv 参考文献条目同样有“内容”和“表现形式”之分，这种可控性是BibTeX带来的。
\end{itemize}


\subsection{参考文献管理器}

参考文献数据库.bib虽然是纯文本的，可以用任意的文本编辑器查看，但总有人喜欢一个找一个``可视化''地查看每一条参考文献。
我想\href{http://jabref.sourceforge.net/}{JabRef}应该是个很不错的选择。
这是一个~Java~写的程序，需要~JRE~才能运行。
就我测试的情况上看，很幸运，JabRef可以顺利打开GBK编码的.bib文件。
但是，打开UTF--8编码的.bib源文件过程中总会崩溃，原因不得而知。
由于我们的~.bib~文件使用的是~UTF-8~编码，所以~JabRef~暂时不可用。

提到参考文献管理器，不得不提到另一个广被使用的软件——\href{http://www.endnote.com/}{EndNote}。
在图书馆的宣讲会上，EndNote被吹得神乎其神，但我发现他对.bib的管理很不友好。
EndNote可以导入.bib文件，却不能导出.bib，只能导出.bbl——被格式化的.bib。
原来，JabRef比较``单纯''，不具备格式化参考文献的能力；
而EndNote有那么一点设置参考文献输出格式的能力，然后就把这种能力滥用，这点搞得我很不爽。
看来，EndNote和Word配合得更好一些。


\section{用~listings~插入源代码}

原先~ctexbook~文档类和~listings~宏包配合使用时，代码在换页时会出现莫名其妙的错误，后来经高人指点，顺利解决了。
感兴趣的话，可以看看\href{http://bbs.ctex.org/viewthread.php?tid=53451}{这里}。
这里给使用~listings~宏包插入源代码的例子，这里是一段C代码。
另外，listings宏包真可谓博大精深，可以实现各种复杂、漂亮的效果，想要进一步学习的同学，可以参考
\href{http://mirror.ctan.org/macros/latex/contrib/listings/listings.pdf}{listings宏包手册}。

\begin{lstlisting}[language={C}, caption={一段C源代码}]
#include <stdio.h>
#include <unistd.h>
#include <sys/types.h>
#include <sys/wait.h>

int main() {
  pid_t pid;

  switch ((pid = fork())) {
  case -1:
    printf("fork failed\n");
    break;
  case 0:
    /* child calls exec */
    execl("/bin/ls", "ls", "-l", (char*)0);
    printf("execl failed\n");
    break;
  default:
    /* parent uses wait to suspend execution until child finishes */
    wait((int*)0);
    printf("is completed\n");
    break;
  }

  return 0;
}
\end{lstlisting}

再给一个插入MATLAB代码的例子，感谢~daisying~站友提供的代码。

\begin{lstlisting}[language={matlab}, caption={一段MATLAB源代码}]
function paper1
r=0.05;
n=100;
T=1;
X=1;
v0=0.8;
sigma=sqrt(0.08);
deltat=T/n;
for i=1:n
    t(i)=i*deltat;
    w(i)=random('norm',0,t(i),1);
end
for i=1:n
    alpha(i)=0.39;
end
for i=1:n
    temp=0;
    for k=1:i
        temp=temp+alpha(k);
    end
    B(i)=exp(r*t(i));
    BB(i)=B(i)*exp(temp*deltat);
    BBB(i)=exp(-r*(T-t(i)));
end
for i=1:n
    s0(i)=X*BBB(i);
    v(i)=v0*exp((r-0.5*sigma^2)*t(i)+sigma*w(i));
    for j=i+1:n
        D=X*BBB(j);
        d1=(log(v(i)/D)+(r+sigma^2/2)*(t(j)-t(i)))/(sigma*sqrt(t(j)-t(i)));
        d2=d1-(sigma*sqrt(t(j)-t(i)));
        ppp(i,j)=D*exp(-r*(t(j)-t(i)))*(1-cdf('normal',d2,0,1))-v(i)*(1-cdf('n
ormal',d1,0,1));
    end
end
for i=1:n
    s1(i)=0;
    for j=i+1:n
        s1(i)=s1(i)+BB(j)^(-1)*alpha(j)*deltat*(X*BBB(j)-B(j)/B(i)*ppp(i,j));
    end
    s2(i)=0;
    for j=1:n
        s2(i)=s2(i)+alpha(j);
    end
    s2(i)=X*exp(-r*T-s2(i)*deltat);
    s(i)=BB(i)*(s1(i)+s2(i));
end
plot(s)
hold on;
plot(s0);
\end{lstlisting}

%%==================================================
%% chapter03.tex for BIT Master Thesis
%% based on CASthesis
%% modified by wei.jianwen@gmail.com
%% modified by 大眼小蚂蚁
%% version: 0.3a
%% Encoding: UTF-8
%% last update: Dec 5th, 2010
%%==================================================

% \bibliographystyle{sjtu2} %[此处用于每章都生产参考文献]

\chapter{问题与希望}
\label{chap:needsomehelp}

由于时间非常仓促，这个模板肯定存在不少问题，所以我需要大家帮助一起解决这
些问题。

下面是我自认为模板需要改进的地方。
\begin{enumerate}
\item \inv 页边距设置上好像有问题，不同打印机打印效果不同；
\item \inv 我只在ubuntu TEXLIVE2011上做过测试，其他环境是否可行，尤为可
  知；
\item \inv 本模板仅仅使用与硕士，不使用与博士、本科；
\item \inv 本文档及其简单
\end{enumerate}

{\heiti 希望有好善者维护BIT论文模板。}



%%==================================================
%% conclusion.tex for SJTU Master Thesis
%% based on CASthesis
%% modified by wei.jianwen@gmail.com
%% version: 0.3a
%% Encoding: UTF-8
%% last update: Dec 5th, 2010
%%==================================================

\chapter*{全文总结\markboth{全文总结}{}}
\addcontentsline{toc}{chapter}{全文总结}

这里是全文总结内容。

 %% 全文总结


%%%%%%%%%%%%%%%%%%%%%%%%%%%%%%
%% 附录(章节编号重新计算，使用字母进行编号)
%%%%%%%%%%%%%%%%%%%%%%%%%%%%%%
\appendix

% 附录中编号形式是"A-1"的样子
\renewcommand\theequation{\Alph{chapter}--\arabic{equation}}
\renewcommand\thefigure{\Alph{chapter}--\arabic{figure}}
\renewcommand\thetable{\Alph{chapter}--\arabic{table}}

%%==================================================
%% app1.tex for SJTU Master Thesis
%% based on CASthesis
%% modified by wei.jianwen@gmail.com
%% version: 0.3a
%% Encoding: UTF-8
%% last update: Dec 5th, 2010
%%==================================================
\chapter{模板更新记录}
\label{chap:updatelog}

\textbf{2012年12月} 根据上海交通大学的模板，按照BIT的规范，构建了此模板。
 % 更新记录
%% app2.tex for SJTU Master Thesis
%% based on CASthesis
%% modified by wei.jianwen@gmail.com
%% version: 0.3a
%% Encoding: UTF-8
%% last update: Dec 5th, 2010
%%==================================================

\chapter{Maxwell Equations}

选择二维情况，有如下的偏振矢量
\begin{subequations}
  \begin{eqnarray}
    {\bf E}&=&E_z(r,\theta)\hat{\bf z} \\
    {\bf H}&=&H_r(r,\theta))\hat{ \bf r}+H_\theta(r,\theta)\hat{\bm
      \theta}
  \end{eqnarray}
\end{subequations}
对上式求旋度
\begin{subequations}
  \begin{eqnarray}
    \nabla\times{\bf E}&=&\frac{1}{r}\frac{\partial E_z}{\partial\theta}{\hat{\bf r}}-\frac{\partial E_z}{\partial r}{\hat{\bm\theta}}\\
    \nabla\times{\bf H}&=&\left[\frac{1}{r}\frac{\partial}{\partial
        r}(rH_\theta)-\frac{1}{r}\frac{\partial
        H_r}{\partial\theta}\right]{\hat{\bf z}}
  \end{eqnarray}
\end{subequations}
因为在柱坐标系下，$\overline{\overline\mu}$是对角的，所以Maxwell方程组中电场$\bf
E$的旋度
\begin{subequations}
  \begin{eqnarray}
    &&\nabla\times{\bf E}=\mathbf{i}\omega{\bf B} \\
    &&\frac{1}{r}\frac{\partial E_z}{\partial\theta}{\hat{\bf
        r}}-\frac{\partial E_z}{\partial
      r}{\hat{\bm\theta}}=\mathbf{i}\omega\mu_rH_r{\hat{\bf r}}+\mathbf{i}\omega\mu_\theta
    H_\theta{\hat{\bm\theta}}
  \end{eqnarray}
\end{subequations}
所以$\bf H$的各个分量可以写为：
\begin{subequations}
  \begin{eqnarray}
    H_r=\frac{1}{\mathbf{i}\omega\mu_r}\frac{1}{r}\frac{\partial
      E_z}{\partial\theta } \\
    H_\theta=-\frac{1}{\mathbf{i}\omega\mu_\theta}\frac{\partial E_z}{\partial r}
  \end{eqnarray}
\end{subequations}
同样地，在柱坐标系下，$\overline{\overline\epsilon}$是对角的，所以Maxwell方程组中磁场$\bf
H$的旋度
\begin{subequations}
  \begin{eqnarray}
    &&\nabla\times{\bf H}=-\mathbf{i}\omega{\bf D}\\
    &&\left[\frac{1}{r}\frac{\partial}{\partial
        r}(rH_\theta)-\frac{1}{r}\frac{\partial
        H_r}{\partial\theta}\right]{\hat{\bf
        z}}=-\mathbf{i}\omega{\overline{\overline\epsilon}}{\bf
      E}=-\mathbf{i}\omega\epsilon_zE_z{\hat{\bf z}} \\
    &&\frac{1}{r}\frac{\partial}{\partial
      r}(rH_\theta)-\frac{1}{r}\frac{\partial
      H_r}{\partial\theta}=-\mathbf{i}\omega\epsilon_zE_z
  \end{eqnarray}
\end{subequations}
由此我们可以得到关于$E_z$的波函数方程：
\begin{eqnarray}
  \frac{1}{\mu_\theta\epsilon_z}\frac{1}{r}\frac{\partial}{\partial r}
  \left(r\frac{\partial E_z}{\partial r}\right)+
  \frac{1}{\mu_r\epsilon_z}\frac{1}{r^2}\frac{\partial^2E_z}{\partial\theta^2}
  +\omega^2 E_z=0
\end{eqnarray}
 % 麦克斯韦方程
% \include{body/app3}


%%%%%%%%%%%%%%%%%%%%%%%%%%%%%%
%% 文后(无章节编号)
%%%%%%%%%%%%%%%%%%%%%%%%%%%%%%
\backmatter

% 参考文献
% 使用 BibTeX
% 包含参考文献文件.bib
{
\zihao{5}
\bibliography{reference/chap1,reference/chap2}
}
% 发表文章目录
%%==================================================
%% pub.tex for SJTU Master Thesis
%% based on CASthesis
%% modified by wei.jianwen@gmail.com
%% version: 0.3a
%% Encoding: UTF-8
%% last update: Dec 5th, 2010
%%==================================================

\begin{publications}{99}

    \item\textsc{Chen H, Chan C~T}. {Acoustic cloaking in three dimensions using acoustic metamaterials}[J].
      Applied Physics Letters, 2007, 91:183518.

    \item\textsc{Chen H, Wu B~I, Zhang B}, et al. {Electromagnetic Wave Interactions with a Metamaterial Cloak}[J].
      Physical Review Letters, 2007, 99(6):63903.
    
\end{publications}


% 致谢
%%==================================================
%% thanks.tex for SJTU Master Thesis
%% based on CASthesis
%% modified by wei.jianwen@gmail.com
%% version: 0.3a
%% Encoding: UTF-8
%% last update: Dec 5th, 2010
%%==================================================

\begin{thanks}

  {\color{red} 希望有牛人来维护咱们学校的模板。}

  感谢母校、感谢联盟！

  感谢高德纳！

  感谢提供交大硕士学位论文~\LaTeX~模板的同学！

  开源精神、奉献精神！

  \vskip 20mm
 {\zihao{-3}
  \hspace{80mm} \lishu{大眼小蚂蚁}

  \hspace{80mm} \lishu{2012年12月}

  \hspace{80mm} \lishu{于北京理工大学9号楼} }
\end{thanks}


\end{document}
